\section{Experimental Methodology}
\label{sec:method}

The experiments are designed to measure $\Crun$ and $\Copt$ for each action under controlled
variation of the workload features, with enough repetition to separate effects from noise.

\subsection{Environment}
All experiments run on one laptop in Spark local mode (Table~\ref{tab:env}). Executors run
inside the driver JVM, so the study measures single-machine, multi-threaded execution. Two
Spark features are disabled to isolate the effects under study \cite{spark2024tuning}. Adaptive Query Execution
(AQE) is off because it coalesces shuffle partitions at run time and would override the A1
setting. Broadcast joins are off so that the join is executed as a sort-merge join and both
inputs are really shuffled. Every requested setting is read back from the running session
and the run is aborted on a mismatch.

\begin{table}[t]
\caption{Study environment}
\label{tab:env}
\centering
\footnotesize
\begin{tabular}{@{}ll@{}}
\toprule
\textbf{Component} & \textbf{Value} \\
\midrule
CPU & Intel i5-1335U, 10 cores (2P+8E), 12 threads \\
Memory / storage & 15.7 GB RAM / NVMe SSD \\
OS & Windows 11 \\
Software & Python 3.11.9, JDK 17, PySpark 3.5.1 \\
Instrumentation & sparkMeasure 0.28 \cite{canali2024sparkmeasure}, psutil \\
Spark mode & \texttt{local[*]}, driver memory 3 GB \\
Fixed settings & AQE off; broadcast join threshold $-1$ \\
\bottomrule
\end{tabular}
\end{table}

\subsection{Workloads}
A seeded generator produces two tables: \texttt{orders} (order id, customer id, product
category, amount, order date) and \texttt{customers} (customer id, region, segment), stored as
Parquet. Three factors are varied. \emph{Data volume} is small (50\,000 orders) or large
(2\,000\,000 orders), with $\max(500, \lfloor\text{orders}/100\rfloor)$ customers.
\emph{Shuffle volume} is controlled through the fraction of orders that pass the pipeline's
filter: about 15\,\% (low) or 85\,\% (high). \emph{Fragmentation} is the number of Parquet
files holding \texttt{orders}: 1 (low), 10 (medium) or 50 (high). Table~\ref{tab:workloads}
lists the ten workloads; all use data seed 42. The A1 workloads fix fragmentation at 10 files
and vary volume and shuffle level; the A2 workloads fix the shuffle level at high and vary
volume and fragmentation. Before the experiments we checked that the factors really change the
amount of data shuffled: at the static setting the median shuffle volume ranges from 0.34\,MB
(small, low) to 34.8\,MB (large, high), and no run spilled to memory or disk.

\begin{table}[t]
\caption{Workloads (data seed 42)}
\label{tab:workloads}
\centering
\footnotesize
\begin{tabular}{@{}lrrrc@{}}
\toprule
\textbf{Workload} & \textbf{Orders} & \textbf{Files} & \textbf{MB} & \textbf{Shuffle} \\
\midrule
A1\_small\_low    & 50\,000      & 10 & 0.59  & low \\
A1\_small\_high   & 50\,000      & 10 & 0.59  & high \\
A1\_large\_low    & 2\,000\,000  & 10 & 20.84 & low \\
A1\_large\_high   & 2\,000\,000  & 10 & 25.31 & high \\
\midrule
A2\_small\_low    & 50\,000      & 1  & 0.55  & high \\
A2\_small\_medium & 50\,000      & 10 & 0.59  & high \\
A2\_small\_high   & 50\,000      & 50 & 0.77  & high \\
A2\_large\_low    & 2\,000\,000  & 1  & 23.31 & high \\
A2\_large\_medium & 2\,000\,000  & 10 & 25.31 & high \\
A2\_large\_high   & 2\,000\,000  & 50 & 27.89 & high \\
\bottomrule
\end{tabular}
\end{table}

\subsection{Pipeline}
The pipeline filters \texttt{orders} on a positive amount, joins the result with
\texttt{customers} on the customer id, and aggregates by customer, region and product
category (sum, count and mean of the amount). The physical plan contains two shuffle
exchanges, one per join input; the aggregation reuses the join's hash partitioning. Each
execution writes its result to Spark's \texttt{noop} sink, which runs the full plan without
producing output files, so that no write cost is mixed into the measurement and no aggregate
column can be pruned. An earlier version of the pipeline, with a broadcast join and a
low-cardinality aggregation, shuffled about 10\,KB regardless of the workload; it was
discarded because no workload factor could influence its shuffle cost.

\subsection{Actions and protocol}
\textbf{A1.} The partition count takes the values $\{1, 2, 4, 8, 16, 50, 200, 400\}$, chosen
from an exploratory sweep to span under-partitioned, near-core-count, default and
over-partitioned settings; 200 is the static baseline. Each of the 32 workload--setting cells
is measured in 5 repetition blocks, giving 160 measured executions. Input tables are cached
once, untimed, before measurement, so that the timed executions reflect computation and
shuffle rather than repeated Parquet reads. Every cell is executed once, untimed, as a session
warm-up. Within each block the cell order is shuffled with a fixed seed (123), and every
measured execution is preceded by one discarded warm-up execution.

\textbf{A2.} Compaction reads the fragmented \texttt{orders} files and rewrites them, using
\texttt{coalesce}, into $\lceil \text{size}/128\,\text{MB}\rceil$ files in a new location; at
our data sizes this yields one file. Coalescing merges partitions without a full shuffle,
which mirrors bin-packing compaction. In every block, each workload is measured in a fixed
sequence: one pipeline execution on the fragmented data (\emph{before}), the compaction job
($\Copt$), and one pipeline execution on the compacted data (\emph{after}). The six workloads
are measured in 5 blocks (90 measured executions), with the workload order shuffled per block
(seed 456). Inputs are \emph{not} cached, because the cost of listing, opening and scanning
many files is exactly what compaction removes. Shuffle partitions stay at the static value of
200, so A2 is measured separately from A1. All three stages of every workload are executed
once, untimed, as a session warm-up; each measured stage is preceded by a discarded warm-up; and
the saving is computed per block from the paired before and after executions.

\textbf{Run hygiene.} Before each session a pre-flight check requires at least 3.5\,GB of free
memory and mains power, and other applications are closed.

\subsection{Metrics}
For every execution we record the wall-clock time of the Spark action (Python
\texttt{perf\_counter}); the stage metrics collected by sparkMeasure
\cite{canali2024sparkmeasure}, including executor run time, executor CPU time, shuffle bytes
read and written, memory and disk spill, garbage-collection time, and the numbers of stages
and tasks; and the system CPU and memory utilization sampled with psutil. Results are appended
to a single CSV file whose header is checked on every write, and every row carries the
pipeline version.

\subsection{Cost estimation and validation}
Per-execution costs are derived from the measurements with proxies P1 and P2
(Section~\ref{sec:problem}) at the EC2 on-demand price of USD 0.0505 per vCPU-hour (m6i
instances, Linux, Asia Pacific (Mumbai) region, accessed 23 October 2026)
\cite{aws2026ec2pricing}. Under P1 the study machine is charged as 12 vCPUs, one per task slot;
the price only converts time into money and does not imply that the laptop performs like an
m6i instance. Every per-execution cost is the median over the repetitions of its cell.
Break-even points are computed for $N \in \{1, 5, 10, 50\}$, and a
sensitivity analysis varies the price by $\pm 50\,\%$ and swaps the proxy. To answer RQ2, the
framework is applied to \tbd{number} unseen validation workloads, generated with different
seeds and an intermediate volume of 500\,000 orders, and its total estimated cost is compared
with that of the static configuration for each $N$.

\subsection{Analysis}
We report medians and interquartile ranges per cell, speed-ups relative to the static setting,
and paired per-block savings for A2. Cells with a coefficient of variation above 0.20 are
flagged and inspected. Outliers are reported, never silently removed; as a sensitivity check
we repeat the A1 analysis with the slowest execution of each flagged cell removed. With five
blocks, an exact two-sided Wilcoxon signed-rank test cannot reach $p < 0.0625$, so we base
conclusions on effect sizes and bootstrap 95\,\% confidence intervals rather than on
significance tests.
